\begin{apartats}

\apartat[1,25]{1}
Aplicant la definició de derivada, calcula la funció derivada de
\[
f(x)=3x^2+x .
\]

\begin{solucio}
\[
f'(x)=\lim_{h\to0}\frac{3(x+h)^2+(x+h)-\left(3x^2+x\right)}{h}
=\lim_{h\to0}\frac{6xh+3h^2+h}{h}=\lim_{h\to0}(6x+3h+1)=6x+1 .
\]
\end{solucio}

\begin{tria}{definicio-tasca}
\itemtria{original}{0,75}{1,25}
Aplicant la definició de derivada, calcula la funció derivada de
\[
g(x)=\sqrt{5-x} .
\]

\begin{solucio}
Multipliquem i dividim pel conjugat:
\[
\begin{aligned}
g'(x)&=\lim_{h\to0}\frac{\sqrt{5-x-h}-\sqrt{5-x}}{h}
=\lim_{h\to0}\frac{(5-x-h)-(5-x)}{h\left(\sqrt{5-x-h}+\sqrt{5-x}\right)}\\
&=\lim_{h\to0}\frac{-1}{\sqrt{5-x-h}+\sqrt{5-x}}=-\frac{1}{2\sqrt{5-x}} ,
\end{aligned}
\]
per a $x<5$. La derivada és negativa: la funció decreix a tot el domini.
\end{solucio}

\itemtria{no-derivable}{0,75}{1,25}
Estudia, aplicant la definició, si la funció
\[
g(x)=|x-3|
\]
és derivable en $x=3$. És contínua en aquest punt?

\begin{solucio}
$g(3)=0$ i $\dfrac{g(3+h)-g(3)}{h}=\dfrac{|h|}{h}$, que val $-1$ si $h<0$ i $1$ si $h>0$.\\
Les derivades laterals són $g'(3^-)=-1$ i $g'(3^+)=1$: són diferents, i per tant $g$
\textbf{no és derivable} en $x=3$ (la gràfica hi fa un pic).\\
En canvi, sí que és contínua: $\lim_{x\to3}|x-3|=0=g(3)$. Una funció contínua en un punt no
hi ha de ser necessàriament derivable.
\end{solucio}
\end{tria}

\begin{nomesllarg}
\apartat{0,75}
Aplicant la definició de derivada, calcula la funció derivada de
\[
k(x)=\frac{1}{x+2} .
\]

\begin{solucio}
\[
k'(x)=\lim_{h\to0}\frac{\frac{1}{x+h+2}-\frac{1}{x+2}}{h}
=\lim_{h\to0}\frac{(x+2)-(x+h+2)}{h(x+h+2)(x+2)}
=\lim_{h\to0}\frac{-1}{(x+h+2)(x+2)}=-\frac{1}{(x+2)^2} ,
\]
per a $x\neq-2$.
\end{solucio}
\end{nomesllarg}

\end{apartats}
